\section{Mathematical notation examples}
\label{app:mathematics}
\draftnote{Template examples only. Adapt the notation and objectives to the
research methods actually used.}

\subsection{Inline and numbered expressions}
Let $\mathcal{D}_{\mathrm{syn}}=\{(\bm{x}_i,y_i)\}_{i=1}^{n}$ denote a
synthetic training dataset, where $\bm{x}_i\in\mathbb{R}^{d}$ is a feature
vector and $y_i$ its target. An illustrative empirical training objective is
\begin{equation}
  \hat{\theta}=\operatorname*{arg\,min}_{\theta\in\Theta}
  \frac{1}{n}\sum_{i=1}^{n}\ell\bigl(f_{\theta}(\bm{x}_i),y_i\bigr).
  \label{eq:training-objective}
\end{equation}
Here $f_{\theta}$ is the model and $\ell$ is its loss function.
Equation~\eqref{eq:training-objective} demonstrates a clickable equation reference.

\subsection{Aligned expressions}
For regression, a squared-error loss and an illustrative real-data evaluation
risk can be written as
\begin{align}
  \ell(\hat{y},y)&=(\hat{y}-y)^2, \label{eq:squared-loss}\\
  R_{\mathrm{real}}(\theta)&=
  \mathbb{E}_{(\bm{x},y)\sim P_{\mathrm{real}}}
  \!\left[\ell\bigl(f_{\theta}(\bm{x}),y\bigr)\right].
  \label{eq:real-risk}
\end{align}
Training on synthetic data does not, by itself, establish performance on real
data; Equation~\eqref{eq:real-risk} defines a separate evaluation quantity.

\subsection{Matrices and cases}
A feature matrix and a binary decision rule illustrate two other useful forms:
\[
  \bm{X}=\begin{bmatrix}
    x_{11} & \cdots & x_{1d}\\
    \vdots & \ddots & \vdots\\
    x_{n1} & \cdots & x_{nd}
  \end{bmatrix},
  \qquad
  \hat{y}=\begin{cases}
    1, & p_{\theta}(y=1\mid\bm{x})\geq\tfrac{1}{2},\\
    0, & \text{otherwise}.
  \end{cases}
\]
